\section*{Chapter \ref{ch:in:portfolio-manager-and-pod-analyst} --- Portfolio Manager and Pod Analyst}

\begin{solution}{exo:in:portfolio-manager-and-pod-analyst:1}
The carried loss is made up first: $40-10=30$ million of net profit, and 15\% of it is \$4.5 million, of which the \$0.5
million salary already paid is an advance. Total for the year: \$4.5 million.
\end{solution}

\begin{solution}{exo:in:portfolio-manager-and-pod-analyst:2}
Expected pay rises by \$7.30 million (\$3.18 to \$10.48 million); the stop probability falls by 42.1 points (57.3\% to
15.2\%).
\end{solution}

\begin{solution}{exo:in:portfolio-manager-and-pod-analyst:3}
About $2\sqrt3=3.46$; with Lo's correction for 756 daily observations, 3.45.
\end{solution}

\begin{solution}{exo:in:portfolio-manager-and-pod-analyst:4}
The payout is a call on each year's profit: the manager shares good years and loses nothing in bad ones but the payouts
that carried losses cancel, so even a zero-mean pod has a positive expected payout. The investors pay for it, as the
netting cost of Book~16, chapter~3; the stop limits it by ending unlucky tenures.
\end{solution}

\begin{solution}{exo:in:portfolio-manager-and-pod-analyst:5}
The certainty equivalent penalises the spread of outcomes: a risk-averse manager values the many low outcomes more than
the few high ones, so she needs more skill before the deal is worth a sure salary.
\end{solution}

\begin{solution}{exo:in:portfolio-manager-and-pod-analyst:6}
The sub-portfolio's volatility and her expected Sharpe ratio (which set the payout's distribution: at 6\% volatility,
\$3 million of annual profit a unit of Sharpe ratio, \$300\,000 at 10\%); the limits and whether a stop ends her job;
whether losses carry forward; what the bonus has been for analysts at her level; and what a record of her own is worth
for her next move.
\end{solution}

\begin{solution}{exo:in:portfolio-manager-and-pod-analyst:7}
At $S=0$ the stop probability falls from 57.3\% to 26.8\% and the expected pay rises from \$3.18 to \$3.56 million; at
$S=1$ the stop probability falls from 15.2\% to 2.1\% and the pay rises from \$10.48 to \$10.76 million. After the cut at
5\% the capital is halved, so a further 5\% of the original capital is a 10\% loss on what is left.
\end{solution}

\begin{solution}{exo:in:portfolio-manager-and-pod-analyst:8}
A Sharpe ratio of 1.2 over three years has a t-statistic of about $1.2\sqrt3=2.08$, barely significant; she was one of
many managers, and the ones who were stopped are not in the sample; 43\% of no-skill tenures survive three years in the
chapter's example. Doubling capital also changes the strategy's costs and capacity. Look at the process, the risk taken
and how concentrated the result was, and raise capital in steps.
\end{solution}

\begin{solution}{pb:in:portfolio-manager-and-pod-analyst:1}
\begin{enumerate}
\item As in the chapter's definition.
\item A percentage of the manager's own profit for the year, mark-to-market, without regard to other managers; losses
carried forward; salary an advance; no clawback; payout due even if the funds lose.
\item It allocates and reallocates capital among managers; capital sets the pod's size, and so its profit and risk.
\item A cut in capital at a first drawdown and a stop at a second; Book~16, chapter~8.
\item The pod's own costs, including its team's pay and data.
\item $\max(w,\rho\max(P_t-L_{t-1},0)K)$ each year with $L_t=\max(L_{t-1}-P_t,0)$.
\item Convex in profit (a call); ended by the stop (a knock-out).
\item \$500 million, 6\% volatility, 15\% payout, \$0.5 million salary, cut at 5\%, stop at 7.5\%, three years.
\item \$3.18 million and 57.3\%; \$10.48 million and 15.2\%; \$23.36 million and 1.9\%.
\item The sure amount with the same expected utility; relative risk aversion 3 and \$5 million of other wealth.
\item 3.85 years of daily results.
\item Capacity, crowding, risk taken, concentration of the result, the process and whether the team moves.
\item A mandate against a benchmark, judged relative to it, with no stop in the deal.
\item Median wages of \$223\,860 for financial managers and \$124\,370 for analysts in the securities industry; they miss
payouts and bonuses.
\item Through a sub-portfolio and a record of her own.
\item $S\approx0.26$, with about 43\% of tenures stopped.
\item With risk aversion 3 and \$5 million of wealth the crossing moves to $S\approx0.75$ (about 23\% stopped).
\item From \$6.04 million and 32.5\% stopped at $S=0.5$ to \$10.48 million and 15.2\% at $S=1$.
\item A guarantee for the first year, the size of the stop, whether capital is cut before the stop, and what happens to
carried losses if she leaves.
\item If her Sharpe ratio is really 0.5 to 1, the deal's expected pay is two to three times her salary even after risk
aversion, but in a third to a sixth of tenures she will be out within three years. Take it only with savings to carry
that outcome and a first-year guarantee.
\end{enumerate}
\end{solution}

\begin{solution}{iq:in:portfolio-manager-and-pod-analyst:1}
Reduce risk before the ladder does it for you, keep the positions with the best expected return per unit of risk, check
whether the loss comes from the process or from one position, and tell the risk manager what you are doing.
\iqlookfor{deliberate de-risking and communication, not hope.}
\end{solution}

\begin{solution}{iq:in:portfolio-manager-and-pod-analyst:2}
With the evidence: when it stopped working, how you know, what it has cost, and what you propose instead, early.
\iqlookfor{honesty with evidence, early.}
\end{solution}

\begin{solution}{iq:in:portfolio-manager-and-pod-analyst:3}
With losses carried forward the payout is far out of the money, and more risk raises its value: an option holder likes
volatility. Platforms answer with risk limits, capital cuts and stops, which cap the risk a manager can add.
\iqlookfor{option convexity and the limits that offset it.}
\end{solution}

\begin{solution}{iq:in:portfolio-manager-and-pod-analyst:4}
Regress the pod's daily returns on market and factor returns; the loadings times the factors' returns is the beta part.
Platforms often want it near zero and may hedge it centrally.
\iqlookfor{factor regression and why the platform cares.}
\end{solution}

\begin{solution}{iq:in:portfolio-manager-and-pod-analyst:5}
Daily returns, the risk and limits under which they were earned, capacity, concentration, turnover, the benchmark and
factor exposures, and whether the team and the data move with the manager.
\iqlookfor{the record's statistics and its conditions.}
\end{solution}

\begin{solution}{iq:in:portfolio-manager-and-pod-analyst:6}
$0.15\,E[\max(X,0)]=0.15\,\sigma/\sqrt{2\pi}=0.15\times30/2.507=\$1.80$ million.
\iqlookfor{the expected positive part of a normal.}
\end{solution}
